% commentaire.tex — f168r, vv. 466–508 (conversion initiale des travaux antérieurs)
\comm{466 sqq.}{}{Scholiaste : \grec{ἐπίτηδες ὁ κηδεστὴς ὑπὲρ τὰς ἄλλας γυναῖκας γυναικίζεται, ἵνα διὰ τούτου λάθῃ}... prudence qui va bientôt le quitter.}
\comm{473}{τί ταῦτ᾽ ἔχουσαι κεῖνον αἰτιώμεθα.}{Ce vers rappelle celui des \italique{Acharniens} (vers 514) :
\citComm{τί ταῦτα τοὺς Λάκωνας αἰτιώμεθα;}
Ils pourraient donc bien être tous deux une reprise parodique du \italique{Télèphe} d'Euripide, pièce maintes fois détournée ici et là.}
\comm{480}{με διεκόρευσεν ἑπτέτιν οὖσαν.}{Mnésiloque n'aura pas mis longtemps à perdre toute retenue : il s'invente bien imprudemment une histoire repoussante qui n'a d'égale que celle de la Quartilla de Pétrone (\italique{Satyrico}, XXV) ou de la marquise de Brinvilliers. (Willems.)}
\comm{484}{στρόφος μ᾽ ἔχει τὴν γαστέρ᾽ ὦνερ κὠδύνη·}{\grec{στρόφος} est le terme «~propre~» pour désigner la \italique{colique} (cf. Hippocrate, \italique{Des maladies aiguës}, XXXVII ; \italique{Aphorismes}, V, 41).}
\comm{486}{κᾆθ᾽ ὁ μὲν ἔτριβεν κεδρίδας ἄννηθον σφάκον·}{Le niais mari se jette sur les graines de genévrier, d'aneth et de sauge car elles sont réputées préserver du mal que sa femme prétend avoir. Voir ce que dit Pline l'Ancien au sujet de ces plantes. (Pline l'Ancien, \italique{Histoire naturelle}. D'après l'édition et la traduction d'Émile Littré, \italique{Histoire naturelle de Pline}, tome second (1877). XXIV, 36 pour le genévrier (ἡ κεδρίς) ; XX, 74 pour l'aneth (τὸ ἄννηθον) ; XXII, 71 pour la sauge (ὁ σφάκος).)}
\comm{487}{ἐγὼ δὲ καταχέασα τοῦ στροφέως ὕδωρ.}{La perfide est rouée (στρόφις) : elle asperge d'eau les gonds de la porte pour ne pas qu'ils fassent de bruit. Le scholiaste l'avait compris ainsi : ἵνα ὀλισθηρὸς γενόμενος μὴ ψοφῇ. On a le même tour chez Plaute (\italique{Curculio}, I, 3, 1) :\citComm{\latin{Anus — Placide egredere, et sonitum prohibe forium, et \textbf{crepitum cardinum},\\
ne quod hic agimus, herus percipiat fieri, mea Planesium.\\
mane, \textbf{subfundam aquolam}.}}J'attire l'attention du lecteur sur le jeu de mots entre καταχέω / καταχέζω d'une part et στροφεύς / στρόφος d'autre part. Je n'en dis pas davantage car en voilà-z-assez !}
\comm{489}{παρὰ τὸν Ἀγυιᾶ κύβδ᾽ ἐχομένη τῆς δάφνης.}{L'Agyieus, surnom tiré d'ἀγυιά, la \italique{rue}, donné à Apollon, considéré comme protecteur des rues, devint par extension le nom des images et autels élevés en l'honneur du dieu, à Athènes, devant les portes des maisons (ἐν προθύροις). (Note d'Émile Saglio\nc{Début de l'article \latin{Agyieus} de l'indispensable \italique{Dictionnaire des antiquités grecques et romaines} de Daremberg \& Saglio ; l'on peut y voir des représentations d'Agyieus.}.) Mnésiloque parle bien d'un autel, ceint selon l'usage de laurier ; Hésychius (α 856) :
\citComm{ἀγυιεύς· ὁ πρὸ τῶν θυρῶν ἑστὼς βωμὸς ἐν σχήματι κίονος\\
\francais{Agyieus : l'autel placé devant les portes, en forme de pilier.}}Sachant ce qu'il y fait, Mnésiloque ajoute donc l'insulte au dieu à l'outrage de l'adultère. Au vu de la forme de l'ἀγυιεύς, on peut imaginer comment il prétend s'être tenu `παρὰ τὸν Ἀγυιᾶ κύβδα', mais pourquoi `ἐχομένη τῆς δάφνης', \italique{s'agrippant au laurier} ?! Quoique porté aux \italique{Realia}, je me suis dispensé de reconstituer la scène ; cela ne m'empêchera pas d'affirmer : voilà une posture fort incommode !
Pourtant, cette image érotique d'un \italique{Penseur} de Rodin (κύβδα) se faisant levrer alors qu'il se tient à un arbre se retrouve ailleurs\nc{Lire Théocrite, \italique{Idylles}, V, 115-119 et voir, par exemple, – ce qui n'est d'ailleurs pas indispensable ! – \italique{Antike Denkmäler}, Band II, 1908, Tafel 41, Figure 3.} ; il y a un lièvre à lever. 
La réponse se trouve peut-être chez les parémiographes : ils ont recensé une des \italique{formules d'Éphèse} (Ἐφέσια γράμματα), incantations magiques pour conjurer le mauvais sort\nc{\italique{Corpus paroemiographorum graecorum}, I, ed. E. L. a Leutsch, Diogeniani centuria IV 78, p. 244. On pourra consulter la \italique{Souda} ou Photius à l'entrée Ἀλεξιφάρμακα : ces charmes pouvaient concerner des époux. Aristophane met dans la bouche de Chrémyle qui prétend avoir quelque succès (\italique{Ploutos}, v. 213) : ὁ Φοῖβος αὐτὸς Πυθικὴν σείσας δάφνην.} :\citComm{Ἐφέσια γράμματα: ἐπῳδαί τινες ἦσαν, ἅσπερ οἱ φωνοῦντες ἐνίκων ἐν παντί. ταὐτὸ τῇ [lire : ταὐτότητι ?]·\
Δαφνίνην φορῶ βακτηρίαν.\\
Lettres éphésiennes : paroles magiques que ceux qui les prononçaient triomphaient en tout. En répétant :
« Je porte la branche de laurier ».}Ainsi, les Grecs \italique{touchaient du bois} (de laurier) pour éviter les malheurs ; cela pourrait expliquer l'attitude de notre libidineuse : s'accrochant aux branches et à de superstitieuses croyances, elle pense ainsi échapper au mauvais œil.}
\comm{494}{ἕωθεν σκόροδα διαμασώμεθα.}{(Pour ne pas que le mari cocu fleure le sacrifice à Vénus de sa femme.) 
À quand une thesmophorienne ecclésie pour châtier Aristophane qui enseigne aux Athéniens toutes les fourberies de leur femme ?.. à moins que le tour ne fût déjà connu (Athénée, XIII, 572c.) :\citComm{Διὰ ταῦθ' ὁ πόρνος οὗτος οὐδὲ τῶν πράσων\\
ἑκάστοτ' ἐπεδείπνει μεθ' ἡμῶν· τοῦτο δ' ἦν,\\
ἵνα μή τι λυπήσειε τὸν ἐραστὴν φιλῶν.}et encore avec de l'oignon (Xénophon, \italique{Symposium}, IV, 8.) :\citComm{Καὶ ὁ Χαρμίδης εἶπεν· `Ὦ ἄνδρες, ὁ Νικήρατος κρομμύων ὄζων ἐπιθυμεῖ οἴκαδε ἐλθεῖν, ἵν᾽ ἡ γυνὴ αὐτοῦ πιστεύῃ μηδὲ διανοηθῆναι μηδένα ἂν φιλῆσαι αὐτόν'.}}
\comm{495}{ἀπὸ τοῦ τείχους εἰσιὼν.}{Thucydide (VIII, 69) :\citComm{ἦσαν [δ᾽] Ἀθηναῖοι πάντες αἰεὶ οἱ μὲν ἐπὶ τείχει, οἱ δ᾽ ἐν τάξει, τῶν ἐν Δεκελείᾳ πολεμίων ἕνεκα ἐφ᾽ ὅπλοις.\\
\francais{Tous les Athéniens restaient continuellement, les uns sur les murailles, les autres à leurs postes, à cause des ennemis qui occupaient Décélie.}}}
\comm{499}{τοὔγκυκλον.}{L'ἔγκυκλον était une sorte de jupe courte, portée par les femmes à l'intérieur. Cf. la note au vers 643.}
\comm{500}{ὑπ᾽ αὐγὰς οἷόν ἐστιν ἐγκεκαλυμμένον.}{Pour le texte, voir le commentaire critique. Le passage, ici comme ailleurs, n'est qu'allusif ; aussi fait-il probablement référence à une histoire bien connue des spectateurs. Victor Coulon\nc{«~Interprétation d'un passage d'Aristophane imité par les conteurs du Moyen-Âge~», in \italique{Annuaire de l'institut de philologie et d'histoire orientales}, tome II (1934) : \italique{Mélanges Bidez}, premier volume, p. 121-34.} l'a rapprochée de contes au motif semblable, couchés par écrit au \nieme{XI} siècle, et de source probablement arabe (parfois, la femme est aidée de sa mère dans son entreprise de mystification du cocu)~; les conteurs orientaux ayant eux-mêmes été familiarisés aux contes grecs. Ainsi, il m'apparaît plus que plausible qu'Aristophane rapporte ici non pas un fait-divers de la chronique scandaleuse du temps mais une historiette populaire ; ces «~révélations~» sur les Athéniennes ne seraient pas si compromettantes.\\Voir le commentaire au vers 643 : l'encycle (ἔγκυκλον) est une courte jupe. Cette donnée finit d'anéantir les hypothèses loufoques où la femme cacherait l'amant aux yeux du mari en tendant l'encycle devant l'amant.}
